\section{The result and its precise scope}\label{sec:scope}

The theorem requested is the strict convex-body theorem:
if \(n\geq1\), \(L=B\Z^n\) is a full-rank lattice, \(K\subseteq\R^n\)
is convex and symmetric about the origin, and
\begin{equation}\label{eq:classical}
   \Vol(K)>2^n|\det B|,
\end{equation}
then \(K\) contains a nonzero point of \(L\).
The familiar proof uses Blichfeldt's principle on \(K/2\), followed by
symmetry and convexity; see the primary exposition \cite{gorodnik}.
The same geometric last step survives here. We replace the continuous
pigeonhole argument by a finite one whose witness is explicitly computed.

There is a necessary language issue. First-order \(\ipra\) quantifies
over natural numbers. It does not literally quantify over arbitrary
subsets of \(\R^n\), real matrices, or Lebesgue measures. An assertion
that a geometric theorem is \emph{in \(\ipra\)} must specify its arithmetic
representation. The following three results do so.

\begin{enumerate}
\item \textbf{Uniform finite arithmetic theorem.}
There is a fixed primitive recursive extractor for rational box
certificates. Its correctness is provable in \(\ipra\). With a rational
basis and a body given by rational symmetric linear inequalities, both
the input checker and the output checker are fixed primitive recursive
functions, uniformly in the dimension.
\item \textbf{Geometric theorem scheme.}
For each arithmetic formula describing membership of rational
lattice-coordinate points, symmetry, midpoint closure, and a finite
inner-volume witness imply a nonzero integer point. Membership need not
be decidable. The arithmetic extractor does not inspect that formula.
\item \textbf{Rational-polytope theorem.}
For bounded rational symmetric polytopes, the usual finite rational
volume inequality produces the required box certificate with an
explicit primitive recursive mesh bound. The weak inequality
\(\Vol(K)\geq2^n|\det B|\) is also sufficient for these closed rational
polytopes.
\end{enumerate}

For represented real bodies, the strict theorem has the same meaning
when positive volume is supplied as effective inner evidence, or by a
representation with a proved procedure that supplies that evidence.
The exact-threshold theorem for arbitrary real bodies is a different
constructive assertion; \cref{sec:obstruction} proves that it entails
weak sign choice.

\subsection{What is proved and what is implemented}
\begin{center}
\begin{tabular}{p{.48\textwidth}p{.43\textwidth}}
\toprule
Deliverable & Status\\
\midrule
Finite certificate theorem and its iPRA induction audit&
Complete human-readable proof below\\
Quantitative rational-polytope reduction and rational equality case&
Complete written argument, including finite volume lemmas\\
Exact rational extraction and witness validation&
Implemented and exercised on the supplied examples\\
Lean data, targets, and completion plan&
Provided as uncompiled source outlines\\
Final Lean term of type \code{Prf PRAAx [] ...}&
Not supplied; an explicit remaining formalization task\\
\bottomrule
\end{tabular}
\end{center}

These distinctions matter. A correct theorem about Lean's \(\R\) or
\(\N\) does not, merely by being accepted by Lean, establish that
\(\ipra\) proves its arithmetic translation. The final formalization
target must assert object-theory derivability.

\section{The base theory and finite coding}\label{sec:arithmetic}

We use first-order intuitionistic logic with equality. Its nonlogical
language contains \(0,S\), and a function symbol for every finite
primitive recursive definition: zero, successor, projections,
composition, and primitive recursion. There are defining equations for
these symbols, successor separation \(Sx\neq0\), successor injectivity,
and induction
\begin{equation}\label{eq:qfind}
 A(0,\vec p)\ \wedge\
 \forall t\bigl(A(t,\vec p)\to A(St,\vec p)\bigr)
 \ \longrightarrow\ \forall t\,A(t,\vec p)
\end{equation}
only when \(A\) is quantifier-free. All parameters \(\vec p\) are allowed.
This is the presentation explicitly defined by \code{PRAAx} in the
inspected \code{ConciseIPRA/PR.lean} \cite{proveit-pr}; compare also
Moschovakis's discussion of the quantifier-free and first-order
presentations \cite{moschovakis}.

Integers are signed natural-number codes. One possible encoding is
\[
 \iota(z)=
 \begin{cases}2z,&z\geq0,\\-2z-1,&z<0.\end{cases}
\]
A rational is an integer numerator and a positive natural denominator.
Internal numerical calculations may use unreduced pairs and compare
them by cross multiplication. For the formula scheme, however, every
rational argument is in \emph{canonical reduced form}: positive
denominator, coprime numerator and denominator, and \(0=0/1\).
All operations feeding that scheme normalize their results.
Reduction is primitive recursive: compute a gcd by bounded divisor
search and divide exactly. The uniqueness of this normal form follows
from the elementary integer divisibility identities in iPRA.
Finite lists use a fixed canonical sequence coding, so equality of
rational vectors gives equality of their codes.
Malformed codes receive designated default decodings, so every function
is total.

\subsection{Bounded statements are Boolean equations}

It is not legitimate to call a formula quantifier-free merely because
all its quantifiers are bounded. Instead, introduce each finite
arithmetic test \(R(i,\vec p)\) by an explicit primitive recursive
characteristic function, or prove in iPRA that \(R\) is equivalent
to that function equaling \(1\). Denote its bit by \([R(i,\vec p)]\).
Then
\[
 \forall i<t\,R(i,\vec p)
 \quad\hbox{is represented by}\quad
 \prod_{i<t}[R(i,\vec p)]=1.
\]
Bounded existentials use finite maxima or a bounded first-witness
function. A finite conjunction of equations and inequalities is likewise
a Boolean equation. All inductions used below are on such equations or
quantifier-free combinations of them.

Natural equality, order, and their finite Boolean combinations are
decidable in this theory. Using a case split on a computed bit therefore
does not introduce unrestricted excluded middle. In particular, a
failed bounded search can be ruled out by arithmetic contradiction,
leaving its successful Boolean branch and its stored witness.

\subsection{The meaning of primitive recursive uniformity}

Our uniform algorithm takes \emph{finite numerical data}: dimensions,
matrices, rational endpoints, and finite lists. No total primitive
recursive interpreter for arbitrary primitive recursive programs is
assumed. Such an interpreter is unavailable. When an application uses
primitive recursive real approximation functions, those functions are
fixed in advance, or their finite computations are supplied as bounded
traces checked by a fixed verifier.

This distinction also prevents a false inference from eventual
termination to availability in \(\ipra\).
Every search in the main theorem has a displayed bound obtained from
its finite input.

\section{Finite inner-volume certificates}\label{sec:certificate}

First work in lattice coordinates. Write
\[
 P=B^{-1}K,\qquad D=\tfrac12P.
\]
Convexity and symmetry are preserved by this linear change of
coordinates, and the desired output is \(z\in\Z^n\setminus\{0\}\)
with \(z\in P\). Its physical lattice vector is \(Bz\).

\begin{definition}[A box certificate]\label{def:cert}
A finite inner-volume certificate for \(D\) is a nonempty finite list
of rational boxes
\[
 Q_j=\prod_{k=0}^{n-1}[\ell_{jk},u_{jk}]
 \quad(0\leq j<s)
\]
such that:
\begin{enumerate}
\item every side length is positive;
\item the interiors of distinct boxes are disjoint;
\item every box is contained in \(D\);
\item their total rational content satisfies
\begin{equation}\label{eq:box-volume}
 W:=\sum_{j<s}\prod_{k<n}(u_{jk}-\ell_{jk})>1.
\end{equation}
\end{enumerate}
\end{definition}

Here containment can mean containment of all rational points; that is
all the arithmetic proof will use. When the boxes are genuinely
contained in a real body, this rational containment follows immediately.
Interiors may be disjoint while boundaries touch. A cell center will be
used so that no choice on a shared boundary is needed.

The certificate gives positive volumetric evidence, not a lattice point:
it consists only of rational endpoints and inclusion facts.
It need not enumerate all points in the body, decide the body's
boundary, or supply the eventual collision.

\subsection{A pure arithmetic checker and a membership scheme}

Let \(\BoxOK(n,c)\) check the numerical parts of
\cref{def:cert}: correct rational codes, positive sides, disjoint
interiors, and \(W>1\). Disjointness is decidable by
\begin{equation}\label{eq:box-disjoint}
 \forall i<j<s\ \exists k<n\
 \bigl(u_{ik}\leq\ell_{jk}\ \vee\
       u_{jk}\leq\ell_{ik}\bigr).
\end{equation}
Every quantifier here is a finite loop compiled into a bit.

Let \(\Phi(q)\) be any fixed arithmetic formula on canonically coded rational
\(n\)-vectors. Additional numerical parameters are suppressed.
Write \(\widehat z\) for the canonical rational embedding of a coded
integer vector \(z\).
Define
\begin{align*}
 \Sym_\Phi&:\quad
   \forall q\,(\Phi(q)\to\Phi(-q)),\\
 \Mid_\Phi&:\quad
   \forall q,r\,\bigl(\Phi(q)\wedge\Phi(r)
       \to\Phi((q+r)/2)\bigr),\\
 \Inside_\Phi(c)&:\quad
   \forall j<s\ \forall q\in Q_j\ \Phi(2q).
\end{align*}
The usual well-formed vector guards are implicit in these displayed
abbreviations and are explicit in a coded version. The interpretation
is \(\Phi(q)\equiv Bq\in K\).

\begin{theorem}[Certificate form of Minkowski in iPRA]\label{thm:main}
There is a fixed primitive recursive function \(F\) such that, for every
formula \(\Phi\) as above, \(\ipra\) proves
\[
\begin{split}
 n\geq1\ \wedge\ \BoxOK(n,c)\ \wedge\
 \Inside_\Phi(c)\ \wedge\ \Sym_\Phi\ \wedge\ \Mid_\Phi\\
 \longrightarrow\
 \bigl(F(n,c)\in\Z^n\ \wedge\ F(n,c)\neq0\
      \wedge\Phi(\widehat{F(n,c)})\bigr).
\end{split}
\]
The output can include an index of a nonzero coordinate and the two
cell indices that justify membership. No induction on \(\Phi\) is used.
Canonical rational operations are essential here: arbitrary formulas
on unreduced fraction codes need not respect equality of rational values.
\end{theorem}

When \(B\) is rational and
\begin{equation}\label{eq:slab}
 K=\{x:\ |a_j\cdot x|\leq b_j\quad(j<t)\}
\end{equation}
has rational rows and bounds, all certificate conditions, including
containment, have one fixed primitive recursive checker.
For a row \(a_j\), put \(\alpha=2a_jB\). Its range on a box is exactly
the interval with endpoints
\[
 \sum_k\min(\alpha_k\ell_k,\alpha_k u_k),
 \qquad
 \sum_k\max(\alpha_k\ell_k,\alpha_k u_k).
\]
Checking that this interval lies in \([-b_j,b_j]\) proves containment.
This yields a single uniform theorem about coded rational data,
as well as the abstract formula scheme.
